\begin{table}[htbp]
\centering
\footnotesize
\caption{截距分解及平台与幂律形式的比较。}
\label{tab:s2}
\begin{tabular}{lllrrrrrr}
\toprule
设定 & $a$ 总量 & $a$ | 周固定效应 & $a$ | 周 + 州季 & 周占比 & ICC & 幂指数 & $p$ & 交叉验证折 \\
\midrule
$R_A$,规则窗口($S=4$) & 0.0150 [0.0100, 0.0208] & 0.0088 [0.0064, 0.0124] & 0.0083 [0.0064, 0.0141] & 0.41 & 0.25 & 0.38 & 0.11 & 26/51 \\
$G$,规则窗口($S=4$) & 0.0840 [0.0508, 0.1226] & 0.0415 [0.0306, 0.0586] & 0.0394 [0.0293, 0.0614] & 0.51 & 0.31 & 0.35 & 0.44 & 19/51 \\
$G$,日历窗口($S=3$) & 0.0991 [0.0583, 0.1399] & 0.0473 [0.0337, 0.0676] & 0.0450 [0.0335, 0.0932] & 0.52 & 0.29 & 0.27 & 0.42 & 28/51 \\
\bottomrule
\end{tabular}
\par\smallskip
\begin{minipage}{0.95\linewidth}\scriptsize 区间为双向自助法 $B=1000$(与表 S1 抽样独立,故略有差异)。周固定效应:拟合前从统计量中去除全国周效应;下一列再去除州-季效应。周占比 = 1 - ($a$|周固定效应)/($a$ 总量)。ICC:统计量方差中周间部分的占比。幂指数:$L=\alpha D^{-c}$ 的 $c$。$p$:仿射平台对幂律的参数自助检验(200 次零分布抽样)。交叉验证折:逐州留一时仿射模型预测更优的折数。功效分析显示,真实截距为估计值的任意比例($\le 1$)时均达不到 80\% 功效。\end{minipage}
\end{table}
